% =====================================================================
% Appendix: proofs for Section 7 (the two-tier learner)
% Sources: research/theory/T7-two-tier-coherent-inferential-learner.md (as repaired in
%          two rounds; see both parts of its "## Verification log"),
%          research/verification/T7-verification.md (incl. "Round 3": independent re-check of
%          the round-2 repairs, no fatal or major findings; its minor fixes are applied),
%          reverification-round2.md, research/theory/T7-checks/*.py (re-run; see app:twotier:scripts).
% =====================================================================
\providecommand{\Steps}{\mathrm{Steps}}
\providecommand{\Ttrust}{T_{\mathrm{trust}}}
\providecommand{\Fres}[1]{F^{(#1)}_{\mathrm{res}}}
\providecommand{\Refut}{\mathrm{Ref}}
\providecommand{\Conf}{\mathrm{Conf}}
\providecommand{\Coll}{\mathrm{Coll}}
\providecommand{\Fsurv}[1]{F^{\mathrm{surv}}_{#1}}
\providecommand{\TTL}{\mathrm{TTL}}
\providecommand{\mgu}{\mathrm{mgu}}
\providecommand{\Prf}{\mathrm{Prf}}
\providecommand{\Sub}{\mathrm{Sub}}
\providecommand{\RobQ}{\mathsf{Q}}
\providecommand{\Ccal}{\mathcal{C}}
\providecommand{\Vcal}{\mathcal{V}}

\section{Proofs for Section~\ref{sec:twotier}}\label{app:twotier}

The proofs of \Cref{prop:twotier:practice}, \Cref{lem:twotier:post} and \Cref{cor:twotier:cpc} are complete in the main text; this appendix adds the details they cite, proves the other results, and contains the examples, computations and secondary results of \cref{sec:twotier} that the main text only points to: the Bayesian remark (\cref{rem:twotier:bayes}), the cost of the audit, the voting audit (\cref{prop:twotier:voting}), the details of the learner's definition, the finite-depth counterexample of \cref{thm:twotier:blame}(c), a Hilbert illustration of collateral, the role of learning when $W$ is a complete checker, and the precise forms of \cref{thm:twotier:arith}(c),(e). Figures marked \status{computed} come from the check scripts listed in \cref{tab:twotier:scripts}. Throughout, (WS) is assumed for every designated position unless a statement says otherwise, $V$ (or $V_{\pos AD}$) denotes a total valuation as in (WS), and ``refutation'' means a refutation in the sense of \Cref{def:twotier:refutation}, with size counting the symbols of distinct judgments. Two facts are used repeatedly.

\begin{itemize}[leftmargin=*]
\item[(F1)] \emph{Cleanness is closed downward.} If $B'\subseteq B$, every refutation of $B'$ is a refutation of $B$, because $R_{B'}\subseteq R_B$. Likewise for step sets.
\item[(F2)] \emph{Transfer.} If every step of a set $S'$ lies in $\Sound(R)$ and $\pi$ is a refutation with steps in $S'\cup R\cup\Ttrust$, then replacing each $S'$-step of $\pi$ by an $R$-derivation of its conclusion from its premises, and compressing (\Cref{def:twotier:refutation}(a)), gives a refutation with steps in $R\cup\Ttrust$, from the same leaves to the same conclusion, of some finite size. In particular, if $R$ is clean at every depth, then every refutation using $S'$ transfers to a refutation of $R$, which is impossible.
\end{itemize}

% ---------------------------------------------------------------------
\subsection{Proof of \texorpdfstring{\cref{lem:twotier:onesided}}{the one-sidedness lemma}}\label{app:twotier:onesided}

Fix a designated position $\pos AD$ and its valuation $V$, and let $X=V^{-1}(1)$. Since $V$ falsifies no step of $\Rstar$, $X$ is closed under $\Rstar$, so $\Cl_{\Rstar}(X)=X$. Hence if $P\step j\in\Sound(\Rstar)$ and $P\subseteq X$, then $j\in\Cl_{\Rstar}(P)\subseteq\Cl_{\Rstar}(X)=X$: \emph{$V$ falsifies no step of $\Sound(\Rstar)$}.

(a) Let $S\subseteq\Sound(\Rstar)$, and let $\pi$ be a derivation relative to $\pos AD$ whose leaves have $e$-value $1$ and whose steps lie in $S\cup\Ttrust$. By induction along the sequence $\pi$, every judgment of $\pi$ is $V$-true. A leaf is $e$-true, hence $V$-true because $V$ extends $e$. A judgment produced by a step of $S\cup\Ttrust$ from earlier, $V$-true judgments is $V$-true, by the claim for $S$ and by (WS) for $\Ttrust$. So the conclusion of $\pi$ is $V$-true, and since $V$ extends $e$ it is not $e$-false. Hence $\pi$ is not a refutation; as $\pi$, its size and the position were arbitrary, $S$ is clean at every depth. For $B\subseteq\Sigma^*$, $R_B\subseteq\Rstar\subseteq\Sound(\Rstar)$ (\Cref{lem:setting:closure}(b)).

(b) A member of $\Conf_d$ has a $d$-refutation, so by (a) it is not contained in $\Sigma^*$; as it is contained in $\Sigma^P=\Sigma^*\sqcup F$, it meets $F$.

(c) The definition of a $d$-refutation of $B$ mentions only $e$ (which is determined by $\Ac$ and $W$), $\Ttrust$, $d$ and $B$. So $\Conf_d$ is a function of these data. Both tie-breaking rules choose among the $d$-refutations of $B$ using only their sizes, a fixed lexicographic order, and whether their descent is blocked; the last depends only on $\pi$ and $e$, because $e^+_\pi$ is computed from $e$ and the trusted steps of $\pi$.

(d) ($\Rightarrow$) The proof of (a) used (WS) only for the position of $\pi$. ($\Leftarrow$) Suppose $e=e_{\pos AD}$ is well defined and $\Sigma^*$ is clean at every depth relative to $\pos AD$. Let $Y=\Cl_{\Rstar\cup\Ttrust}(A\cup W^{-1}(1))$ and let $V$ be the indicator of $Y$. Since $Y$ is closed under $\Rstar\cup\Ttrust$, $V$ falsifies no step of $\Rstar\cup\Ttrust$, and $V$ is $1$ on every $e$-true judgment. Let $j$ be $e$-false and suppose $j\in Y$. By \Cref{lem:setting:closure}(a), $j$ has a finite $(\Rstar\cup\Ttrust)$-derivation tree from $A\cup W^{-1}(1)$, that is, from $e$-true leaves. Compressing it gives a refutation of $\Sigma^*$ relative to $\pos AD$ of some finite size, contradicting cleanness. So $V$ is $0$ on every $e$-false judgment, $V$ extends $e$, and (WS) holds for $\pos AD$. Nothing in the argument used that $\Sigma^*$ is the target. \qed

% ---------------------------------------------------------------------
\subsection{Details for \texorpdfstring{\cref{prop:twotier:practice}}{the practice proposition}}\label{app:twotier:practice}

The first two claims are \Cref{thm:imitation:iidnoise,thm:imitation:trimmed} applied to the tagged practice, together with \Cref{def:setting:practice}, which gives every fallacy an instance outside $\Sound(\Rstar)$. For the $\CPC$ claim, let $\Pi_c\step j_c$ be the closed falsified instance of the pure fallacy $\tau$ (\Cref{lem:twotier:post}). Each closed formula has the same value under every valuation, so every member of $\Pi_c$ and $\neg j_c$ are tautologies, and a complete $\Sigma^*$ derives them from $\emptyset$. The step adds $j_c$, and $\{j_c,\neg j_c\}$ classically entails every formula, which $\Sigma^*$ derives. For the contrast: if $\Sigma^*$ is complete and $s=(p\step q)$ for atoms $p,q$, the set of tautologies is closed under $s$ because $p$ is not a tautology, so $\Cl_{\Rstar\cup\{s\}}(\emptyset)=\Cl_{\Rstar}(\emptyset)$. \qed

% ---------------------------------------------------------------------
\subsection{Proof of \texorpdfstring{\cref{lem:twotier:reiter}}{the hitting-set duality}}\label{app:twotier:reiter}

(a) $M\subseteq U$ is clean iff it contains no member of $\mathcal K$, iff (by upward closure and finiteness) it contains no member of $\Ccal$, iff $U\setminus M$ meets every member of $\Ccal$, i.e.\ is a transversal. Complementation reverses inclusion, so maximal clean sets correspond to minimal transversals.

(b, $\Leftarrow$) Let $x\in C\in\Ccal$ and $T=(U\setminus C)\cup\{x\}$. Let $E\in\Ccal$. If $E\not\subseteq C$, then $E$ meets $U\setminus C$. If $E\subseteq C$, then $E=C$ by minimality of $C$, so $x\in E$. Hence $T$ is a transversal, and it contains a minimal transversal $T'$. Since $T'$ meets $C$ and $T'\cap C\subseteq T\cap C=\{x\}$, we get $x\in T'$.

(b, $\Rightarrow$) If $x\in T$ for a minimal transversal $T$, then $T\setminus\{x\}$ is not a transversal, so some $E\in\Ccal$ has $E\cap T=\{x\}$, and $x\in E$.

(b, second claim) By (a), $x$ lies outside some maximal clean set iff $x$ lies in some minimal transversal, iff $x\in\bigcup\Ccal$. Every clean set is contained in a maximal one, and $\emptyset$ is clean, so there are maximal clean sets and their intersection is $U\setminus\bigcup\Ccal$. It is clean because cleanness is closed downward.

(c, $\Leftarrow$) If $\Ccal=\{\{x_1\},\dots,\{x_r\}\}$, every transversal contains $\{x_1,\dots,x_r\}$, which is itself a transversal.

(c, $\Rightarrow$) Let $T$ be the unique minimal transversal. By (b), $T=\bigcup\Ccal$. For each $x\in T$, minimality gives some $E\in\Ccal$ with $E\cap T=\{x\}$; since $E\subseteq\bigcup\Ccal=T$, $E=\{x\}$. A member $C\in\Ccal$ with $|C|\ge2$ would then strictly contain the singleton member $\{x\}$ for any $x\in C$, contradicting minimality. \qed

{\sloppy The Lean formalization (\texttt{InfLearn.Blame}) proves (a) as \texttt{isMaxClean\_iff} and \texttt{setOf\_isMaxClean\_eq\_image}, (b) as \texttt{exists\_minTransversal\_mem\_iff}, \texttt{iInter\_maxClean\_eq} and \texttt{isClean\_sdiff\_blameSet}, and (c) as \texttt{existsUnique\_minTransversal\_iff}, for a finite universe; upward closure is needed only inside $U$, and $\emptyset\notin\mathcal K$ only for the intersection claims.\par}

% ---------------------------------------------------------------------
\subsection{Proof of \texorpdfstring{\cref{prop:twotier:dilemma}}{the dilemma}}\label{app:twotier:dilemma}

(a) $\emptyset\subseteq\Sigma^*$ is $d$-clean by \Cref{lem:twotier:onesided}(a), so $\emptyset\in\Vcal_d$ and $\bigcap_{\Sigma\in\Vcal_d}R_\Sigma\subseteq R_\emptyset=\emptyset$. The version-space verifier (\Cref{def:caution:vs}) accepts exactly this intersection.

(b) Before any refutation is found, $\Vcal'_t=\Vcal$, whose unique maximal member is $\Sigma^P$; so the verifier accepts $R^P$, which contains an instance outside $\Sound(\Rstar)$ of every $\tau\in F$.

(c) \emph{Schema level.} $\Vcal_d$ is the family of clean sets for $U=\Sigma^P$ and $\mathcal K=\Conf_d$, which is upward closed and does not contain $\emptyset$. By \Cref{lem:twotier:reiter}(b) the intersection of its maximal members is $\Sigma^P\setminus\bigcup\Ccal_d$, and since $\Sigma^P=\Sigma^*\sqcup F$ this equals $(\Sigma^*\setminus\Coll_d)\cup\Fsurv d$. For the inclusion $\Fsurv d\subseteq\Fres d$, let $\tau\in F\setminus\Fres d$. Then $\Sigma^*\cup\{\tau\}\in\Conf_d$ contains some $C'\in\Ccal_d$, and $C'\not\subseteq\Sigma^*$ by \Cref{lem:twotier:onesided}(a), so $\tau\in C'$ and $\tau\notin\Fsurv d$. Hence the asserted set lies in $\Rstar\cup R_{\Fres d}$.

\emph{Instance level.} Let $\mathcal M$ be the maximal members of $\Vcal_d$. For each $\Sigma\in\mathcal M$, $\bigcap\mathcal M\subseteq\Sigma$ gives $R_{\bigcap\mathcal M}\subseteq R_\Sigma$; so the instance-level set contains the schema-level one. Since $\Sigma^*\in\Vcal_d$ and $\Sigma^P$ is finite, $\Sigma^*\subseteq\Sigma_{\max}$ for some $\Sigma_{\max}\in\mathcal M$. For $\tau\in\Sigma_{\max}\cap F$, $\Sigma^*\cup\{\tau\}\subseteq\Sigma_{\max}$ is $d$-clean by (F1), so $\tau\in\Fres d$. Hence $\bigcap_{\Sigma\in\mathcal M}R_\Sigma\subseteq R_{\Sigma_{\max}}\subseteq\Rstar\cup R_{\Fres d}$. As a step set it is $d$-clean, since a $d$-refutation of it would be a $d$-refutation of the schema set $\Sigma_{\max}$. The MP/AC example of the main text shows that the inclusion can be strict: $(A,\ A\to A\step A)$ is the instance of MP with $B:=A$ and the instance of AC with $B:=A$, so its instances lie in $R_{\mathrm{MP}}\cap R_{\mathrm{AC}}$, while $R_{\{\mathrm{MP}\}\cap\{\mathrm{AC}\}}=\emptyset$. \qed

\begin{remark}[A Bayesian middle way]\src{T7 Remark 2.4$'$}\ \status{proof sketch}\label{rem:twotier:bayes}
Put a prior $w$ on $\Vcal$ and accept $s$ iff $w\{\Sigma\in\Vcal'_t:s\notin R_\Sigma\}<w_0$. By the argument of \Cref{thm:caution:bayesdet} this is sound at all times for every target with $w(\Sigma^*\cup\Fres d)\ge w_0$, provided that set is clean, and it waits for the audit automatically. It cannot break the symmetry of \Cref{thm:twotier:blame}, where the rival diagnosis has the same prior mass. We do not develop it; the audit-then-assert design of TTL is simpler and has explicit constants.
\end{remark}

% ---------------------------------------------------------------------
\subsection{Proof of \texorpdfstring{\cref{lem:twotier:descent}}{the descent lemma}}\label{app:twotier:descent}

\emph{Agreement.} We show by induction on the order in which the fixed-point computation assigns values that every value of $e^+_\pi$ agrees with $V$. Initial values come from $e$, which $V$ extends. If the forward rule assigns $1$ to the conclusion of a trusted step whose premises have value $1$, those premises are $V$-true by induction, and since $V$ falsifies no trusted step the conclusion is $V$-true. If the backward rule assigns $0$ to a premise $\pi_0$ of a trusted step whose conclusion has value $0$ and whose other premises have value $1$, then by induction the conclusion is $V$-false and the other premises are $V$-true; if $\pi_0$ were $V$-true the step would be falsified by $V$, so $\pi_0$ is $V$-false. Hence $e^+_\pi$ never assigns both values to a judgment, and if it does, no total $V$ as in (WS) exists for $\pos AD$.

(a) Each move of the descent goes from a node of value $0$ to one of its premises, which occurs strictly earlier in the sequence $\pi$; so the descent terminates. Leaves have $e$-value $1$, hence $e^+_\pi$-value $1$, while every visited node has value $0$; so every visited node is the conclusion of some step of $\pi$. Since the descent is unblocked, it ends by outputting a step $s$ all of whose premises have value $1$ and whose conclusion has value $0$ (a step without premises is output at once). By agreement, the premises of $s$ are $V$-true and its conclusion $V$-false. By (WS), $s\notin\Ttrust$ and $s\notin\Rstar$; as $\pi$ is a refutation of $B$, $s\in R_B$.

(b) The worklist assigns each judgment of $\pi$ a value at most once and re-examines a trusted step only when one of its judgments receives a value, so the fixed point costs $O(|\pi|\varphi)$ operations, with one evaluation of $e$ per judgment. The descent looks at the premises of each visited node once, and visits at most $|\pi|$ nodes. The number of judgments is at most $|\pi|\le d$.

(c) If some genuine $\sigma$ had $s\in\inst(\sigma)$, then $s\in\Rstar$, contradicting (a).

(d) Let $\sigma\in B$ with $s\in\inst(\sigma)$. Let $J_1$ be the list of judgments of $\pi$ with $e^+_\pi$-value $1$, in the order in which they received the value (those with initial $e$-value $1$ first). Each member of $J_1$ is $e$-true, and then is treated as a leaf, or received value $1$ by the forward rule as the conclusion of a trusted step of $\pi$ whose premises precede it in $J_1$. The premises of $s$ lie in $J_1$. Let $c_0$ be the conclusion of $s$. Define $c_0,c_1,\dots$: if $c_i$ is $e$-false, stop; otherwise $c_i$ received value $0$ by the backward rule from a trusted step $t_i$ of $\pi$ of which $c_i$ is a premise, whose conclusion $c_{i+1}$ had received value $0$ strictly earlier, and whose other premises lie in $J_1$. Since assignment times strictly decrease along the chain, it terminates at some $e$-false $c_k$, and the $c_i$ are pairwise distinct. They are not in $J_1$, because $e^+_\pi$ is consistent.

Let $\pi_\sigma$ be the list $J_1$, followed by $c_0$ (justified by $s$), $c_1$ (justified by $t_0$), \dots, $c_k$ (justified by $t_{k-1}$). Its judgments are distinct; each non-leaf judgment is the conclusion of a step whose premises occur earlier; its leaves are $e$-true; its steps are $s\in\inst(\sigma)$ and trusted steps; and its last judgment $c_k$ is $e$-false. So $\pi_\sigma$ is a refutation of $\{\sigma\}$ relative to $\pos AD$. All its judgments are judgments of $\pi$, and size counts distinct judgments, so $|\pi_\sigma|\le|\pi|\le d$. Hence $\{\sigma\}$ is a $d$-conflict. (Under a tree-size convention this bound could fail, because the premise derivations of $s$ and the side premises of the chain may share judgments.) \qed

The Lean theorem \texttt{InfLearn.Blame.LineArg.descent} proves the descent for a line-list argument under a total evaluation, with the walk length bounded by any height function that decreases along dependencies, and \texttt{LineArg.descent\_not\_mem} proves that the step found lies outside every rule set whose steps preserve truth. Trusted-step propagation and part (d) are not formalized.

% ---------------------------------------------------------------------
\subsection{Proof of \texorpdfstring{\cref{lem:twotier:audit}}{the audit lemma}, and the cost of the audit}\label{app:twotier:audit}

By \Cref{lem:twotier:descent}(a),(c) each successful descent removes at least one schema, all of them fallacies, so $\Sigma^*\subseteq B$ throughout, and there are at most $|F|$ removals. (a) The returned $B$ is $d$-clean; for $\tau\in B\setminus\Sigma^*$, $\Sigma^*\cup\{\tau\}\subseteq B$ is clean by (F1), so $\tau\in\Fres d$. (b) By \Cref{lem:twotier:reiter} with $U=B_0$, $A$ is the intersection of the maximal clean subsets of $B_0$, hence clean. The left inclusion holds because $\Sigma^*\subseteq B_0$ and $\Coll_d(B_0)=\Sigma^*\cap\bigcup\Ccal_d(B_0)$. If $\tau\in A\cap F$ were not residual, $\Sigma^*\cup\{\tau\}\subseteq B_0$ would contain a minimal conflict, which contains $\tau$ by \Cref{lem:twotier:onesided}(a), contradicting $\tau\in A$. The last claim is \Cref{prop:twotier:dilemma}(c). \qed

\emph{Cost.} The descent path costs at most $|F|+1$ oracle calls and $O((|F|+1)\,d\,\varphi)$ evaluations of $e$. The fallback needs $\bigcup\Ccal_d(B_0)$: at most $2^{|B_0|}$ oracle calls by brute force. From a cleanness (membership) oracle, the minimal conflicts and the maximal clean sets of a monotone property can be generated jointly in time quasi-polynomial in their total number \citep{bioch1995complexity,gurvich1999generating}, with the dualization algorithm of \citet{fredman1996complexity} as the engine; the second family can be exponentially larger than the first, and both can be exponential in $|F|$. At stable depth the fallback's schema-level incompleteness is unavoidable (\Cref{thm:twotier:blame}(b)); its cost is the price of bags.

% ---------------------------------------------------------------------
\subsection{The voting audit}\label{app:twotier:voting}

\Cref{thm:twotier:main} assumes (WS) for \emph{every} designated position. When some positions may be wrong, as when an idealization is designated together with background it contradicts, the audit can vote.

\begin{proposition}[Voting audit]\src{T7 Prop 3.2, as repaired}\label{prop:twotier:voting}
Suppose at most $m$ designated positions violate (WS), it is not known which, and $\hat P=\Sigma^P$. For each position $p$ run the descent loop of the audit relative to $p$ alone, from $B_p=\Sigma^P$ and $\mathrm{Blamed}_p=\emptyset$, with the prefer-unblocked oracle restricted to $p$. Stop when the oracle returns NONE ($p$ is \emph{finished}), when the descent is blocked, or when $e_p$ is ill defined or $e^+_\pi$ gives a judgment both values (then $p$ is certified bad and its blame is discarded); otherwise move every $\sigma\in B_p$ with $s\in\inst(\sigma)$ into $\mathrm{Blamed}_p$. Let $\mathrm{votes}(\sigma)=\#\{p:\sigma\in\mathrm{Blamed}_p\}$, counting blame from $\pos\emptyset\emptyset$ as $m+1$ votes if the world is known to be truthful, and return $A=\Sigma^P\setminus\{\sigma:\mathrm{votes}(\sigma)\ge m+1\}$. Then
\textup{(a)} $\Sigma^*\subseteq A$;
\textup{(b)} at most $|\Ac|(|F|+1)+m|\Sigma^*|$ oracle calls are made;
\textup{(c)} call $p$ \emph{complete} if it is good and finished; then $A\cap F\subseteq\{\tau\in F:\#\{p\text{ complete}:\Sigma^*\cup\{\tau\}\text{ is not $d$-clean relative to }p\}\le m\}$. For $m=0$ and all positions complete this is \Cref{lem:twotier:audit}(a).
\end{proposition}

\begin{proof}[Proof]
(a) At a good position $p$, every schema moved into $\mathrm{Blamed}_p$ is a fallacy, by \Cref{lem:twotier:descent}(c) applied relative to $p$. So a genuine schema collects votes only from bad positions, of which there are at most $m$. The bonus for $\pos\emptyset\emptyset$ is used only when that position is known to be good, and then it blames only fallacies. So every genuine schema has at most $m$ votes and lies in $A$.

(b) Every call at $p$ either ends $p$'s loop or moves at least one schema out of $B_p$. Indeed, if the call returns a refutation $r$, $e_p$ is well defined, $e^+_{p,r}$ is consistent and the descent is unblocked, then the descent outputs a step $s$ whose premises have value $1$ and whose conclusion has value $0$, by the argument of \Cref{lem:twotier:descent}(a), which used only consistency and unblockedness. If $s$ were trusted, the forward rule would have given its conclusion value $1$ as well, contradicting consistency. So $s$ is untrusted and lies in $R_{B_p}$, and some $\sigma\in B_p$ with $s\in\inst(\sigma)$ is moved. Hence the loop at $p$ makes at most $|B_p|+1\le|\Sigma^P|+1$ calls, and at most $|F|+1$ at a good position, where only fallacies move. With $m'\le m$ bad positions the total is at most $(|\Ac|-m')(|F|+1)+m'(|F|+|\Sigma^*|+1)\le|\Ac|(|F|+1)+m|\Sigma^*|$.

(c) Let $p$ be complete. Its loop ended with NONE, so its final $B_p$ is $d$-clean relative to $p$; and since $p$ is good, only fallacies were moved, so $\Sigma^*\subseteq B_p$. If $\tau\in F\setminus\mathrm{Blamed}_p$, then $\tau\in B_p$, and $\Sigma^*\cup\{\tau\}\subseteq B_p$ is $d$-clean relative to $p$ by (F1). Contrapositively, each complete $p$ relative to which $\Sigma^*\cup\{\tau\}$ is not $d$-clean gives $\tau$ a vote. A $\tau\in A$ has at most $m$ votes, which gives the claim. For $m=0$ with all positions complete, $A=\Sigma^P\setminus\bigcup_p\mathrm{Blamed}_p$ lies inside every final $B_p$, so it is $d$-clean relative to every position; it contains $\Sigma^*$, and each of its fallacies is residual. This is the conclusion of \Cref{lem:twotier:audit}(a).
\end{proof}

\emph{Why per-position loops} \status{computed}. One smallest-refutation query per position, followed by voting, does not suffice. With one bad position $\pos{a,\ a\to b}{b}$, target $\{\mathrm{MP}\}$ and $F=\emptyset$, a single query already makes a successful descent, which blames the genuine MP; so the number of successful descents is not bounded by $(m+1)|F|$, and votes from bad positions must be tolerated, as in (a) and (b). With $m=1$, the truthful positions $p_1=\pos{q,\ p\to q,\ \neg p}{p,\neg q}$ and $p_2=\pos{\neg r,\ r\to t}{\neg t}$ and the practice $\{\mathrm{MP},\mathrm{AC},\mathrm{DA}\}$, the smallest refutation at $p_1$ always blames AC and hides denying the antecedent, so DA gets one vote and survives although it is refutable at two positions. The loops above exclude already-blamed schemas, so DA is blamed at both positions and removed, while AC, refutable at one position only, may survive.

% ---------------------------------------------------------------------
\subsection{Proof of \texorpdfstring{\cref{prop:twotier:sandbox}}{the sandbox proposition}}\label{app:twotier:sandbox}

(a) By \Cref{def:twotier:ttl}, $A$ is computed at $t=N_1$ from $\hat P$, a function of the first $N_1$ human steps, and from answers of $\Refut_d$, which are deterministic functions of $(B,d,\Ac,W,\Ttrust)$ (\Cref{lem:twotier:onesided}(c)). Phase-1 acceptance is membership in $R_A$. The sandbox, the red team and the prover interact with the tier only through queries, which do not change $A$, and escalation answers are not added to the sample. A promoted $\kappa$ is a macro, answered by expanding it into an $R_A$-derivation, so the primitive accepted steps remain $R_A$ and $\Cl_{R_A}$ is unchanged. If promoted rules were accepted as primitive steps instead, then $\kappa\subseteq\Sound(R_A)\subseteq\Sound(\Rstar\cup R_{\Fres d})$, so the accepted set would satisfy only this weaker inclusion; the closure would be unchanged (\Cref{lem:setting:closure}(b)). But $d$-cleanness of the enlarged step set can fail. If some $\tau\in A$ is residual at depth $d$ with a shortest refutation of size $2d$, a promoted $\kappa$ whose instances package the first half of that refutation into single steps removes its intermediate judgments and yields a refutation of size at most $d$. Primitive promoted rules would also need truth maintenance under the depth schedule.

(b) Let $h_0\in\Hc_S$ be clean at every depth with $w(h_0)>0$. A detection exhibits a refutation $\pi$ of the step set $B_t=\bigcap S_t$ and deletes every $h\supseteq\Steps(\pi)\setminus\Ttrust$. If $h_0$ were deleted, the steps of $\pi$ would lie in $h_0\cup\Ttrust$ and $\pi$ would be a refutation of $h_0$, which is impossible. So $h_0$ is never deleted. In a detection, every $h\in S_t$ contains $B_t\supseteq\Steps(\pi)\setminus\Ttrust$, so all of $S_t$ is deleted and $w(\VS_{t+1})\le\frac12w(\VS_t)$. After $D$ detections $w(h_0)\le w(\VS)\le2^{-D}$, so $D\le\log_2(1/w(h_0))$. For the robust bound, run the proof of \Cref{thm:coherence:robust} with $h_0$ in place of $\Rstar$: every detection multiplies by $\beta$ the weight of a superset of $S_t$, so the total weight shrinks by a factor at most $\frac{1+\beta}2$; and $h_0$ is penalized only in false alarms, since a refutation of $h_0$ can only come from a position violating (WS). Hence $\beta^mw(h_0)\le(\frac{1+\beta}2)^D$. Deletion must test only the \emph{untrusted} steps of $\pi$: if it tested all steps against hypotheses that do not contain $\Ttrust$, coalition members whose refutations use trusted steps would never be deleted, and halving would fail.

(c) For finite $S_t$, $s\in\bigcap_{h\in S_t}h$ iff for each $h$ (a finite set of schemas) $s$ matches some schema of $h$, which is decidable. For schema-valued hypotheses, $\bigcap_h\inst(\sigma_h)$ is empty or the instance set of the most general common instance, obtained by unifying renamed-apart copies \citep{robinson1965machine}. \qed

% ---------------------------------------------------------------------
\subsection{Proof of \texorpdfstring{\cref{thm:twotier:main,cor:twotier:truth}}{the main theorem}}\label{app:twotier:main}

\emph{Step 1: identification (noisy mode).} For each tag $i$ consider the following events on the first $N_1$ samples: (I$_i$) at most $e_i$ invalid steps are tagged $i$; and (II$_i$) the valid tag-$i$ samples are $e_i$-robustly generic (\Cref{thm:imitation:trimmed}(b)). As in the proof of \Cref{thm:imitation:iidnoise}, (II$_i$) follows from $c_i$ count events: for each metavariable $x$ with best split $G_{i,x}$ (achieving $\rho_{i,x}$), more than $e_i$ valid tag-$i$ samples have root in $G_{i,x}$ and more than $e_i$ have root outside it; and for each pair $x\neq y$, more than $e_i$ have $\Theta x\neq\Theta y$. (For any $f$, the samples on the side of $G_{i,x}$ not containing $f$ all have root $\neq f$.) For a ground tag the single event is that more than $e_i$ valid samples occur; this is the convention $\rho_i=c_i=1$. Let $G$ be the intersection of all these events.

The number of invalid tag-$i$ steps is $\mathrm{Bin}(N_1,\alpha_i)$ with $\alpha_i\le\bar\alpha_i$. Since $e_i=\lfloor(\bar\alpha_i+\Delta)N_1\rfloor$, exceeding $e_i$ means exceeding the mean by more than $\Delta N_1$, which by Hoeffding's inequality \citep{hoeffding1963probability} has probability at most $e^{-2N_1\Delta^2}$. Each count in (II$_i$) is $\mathrm{Bin}(N_1,p)$ with $p\ge\rho_i\beta_i\ge\bar\alpha_i+2\Delta$ by (Floor); being at most $e_i\le(\bar\alpha_i+\Delta)N_1$ means falling short of the mean by at least $\Delta N_1$, again with probability at most $e^{-2N_1\Delta^2}$. A union bound over the $1+c_i$ events per tag gives
\[
\Prob(G^c)\le\sum_i(1+c_i)e^{-2N_1\Delta^2}\le(1+c)Ke^{-2N_1\Delta^2}\le\delta
\]
by the choice of $N_1$. On $G$, \Cref{thm:imitation:trimmed}(b), applied to the realizable tagged practice, shows that the trimmed verifier of tag $i$ accepts exactly $\inst(\sigma_i)$.

It remains to see that $\hat\sigma_i$ computes this set. On $G$, $n_i>e_i$. The trimmed version space of tag $i$ consists of the hypotheses $\inst(\tau)$ (and $\emptyset$) with $|P_i\setminus\inst(\tau)|\le e_i$; $\emptyset$ is excluded because $n_i>e_i$. Each $\inst(\lgg(P_i\setminus E))$ with $|E|=e_i$ belongs to it. Conversely, if $|P_i\setminus\inst(\tau)|\le e_i$, then $P_i\setminus E\subseteq\inst(\tau)$ for some $E$ of size $e_i$, so $\tau\succeq\lgg(P_i\setminus E)$ (\Cref{thm:setting:lgg}). Hence the accepted set is $\bigcap_{|E|=e_i}\inst(\lgg(P_i\setminus E))$, a finite intersection of instance sets. It equals $\inst(\sigma_i)\neq\emptyset$, and a finite intersection of instance sets with a common instance is the instance set of the most general common instance. So $\inst(\hat\sigma_i)=\inst(\sigma_i)$, hence $\hat\sigma_i\equiv\sigma_i$, and $\hat P=\Sigma^P$ up to renaming.

\emph{Step 1 (noise-free mode).} Now $e_i=0$ and $\hat\sigma_i=\lgg(P_i)$. Let $G$ be the event that every tag's samples satisfy the witness events (R) and (D) of \Cref{lem:setting:recover} (for a ground tag: that it is sampled). By \Cref{thm:imitation:coupon}, with its convention $\rho_i=c_i=1$ for ground rules, $\Prob(G^c)\le\sum_ic_ie^{-N_1\pi_i\rho_i}\le cKe^{-N_1\lambda}\le\delta$, and on $G$, $\lgg(P_i)\equiv\sigma_i$. In both modes $G$ depends only on the first $N_1$ human samples; human answers to escalated queries are not part of the sample.

\emph{Two details of the definition} \status{computed}. In noise-free mode the budget must be $e_i=0$: with $\bar\alpha_i=0$ the noisy formula still gives $e_i=\lfloor\Delta N_1\rfloor>0$, and then the event of \Cref{thm:imitation:coupon} does not imply that trimmed lggs identify the rule. With one tag, one metavariable, two equiprobable roots ($\lambda=\frac12$) and $\delta=0.001$, $N_1=16$; identification fails with probability $0.077$ for $e=4$ and $3\cdot10^{-5}$ for $e=0$. And with $n_i\le e_i$ samples the trimmed version space contains the empty rule, so its intersection is $\emptyset$, not ``everything''; coding it as ``accept all'' made TTL unsound in 12 of 20 simulated runs at $N=60$.

\emph{Step 2: audit.} On $G$, $\hat P=\Sigma^P$, and \Cref{lem:twotier:audit} applies. It gives the two cases of (iii), the $d$-cleanness of $A$ in (i), and the audit part of (ii), with the evaluation count from \Cref{lem:twotier:descent}(b). For the completeness claim in (iii) under (SB$_d$): every member of $\Ccal_d$ is a singleton $\{\tau\}$, and $\tau\in F$ by \Cref{lem:twotier:onesided}(b), so $\Coll_d=\emptyset$; and $\Coll_d(B_0)\subseteq\Coll_d$, because a minimal member of $\Conf_d\cap2^{B_0}$ is a minimal member of $\Conf_d$ (its proper subsets lie in $2^{B_0}$).

\emph{Step 3: soundness.} Before $N_1$ the tier accepts nothing. From $N_1$ on it accepts exactly $R_A$, and $A\subseteq\Sigma^*\cup\Fres d$ gives $R_A\subseteq\Rstar\cup R_{\Fres d}=:R'$. Since $R'\subseteq\Sound(R')$, \Cref{lem:search:stepwise} with $R'$ as target gives $\Cl_{R_A}(B)\subseteq\Cl_{R'}(B)$ for all $B$; the version with trusted steps follows from monotonicity of $\Cl$ in the rule set (\Cref{lem:setting:closure}(a)). If $\Fres d=\emptyset$ then $R'=\Rstar$. The accepted set does not depend on the prover's queries, so the conclusions hold for every prover.

\emph{Step 4: sandbox.} \Cref{prop:twotier:sandbox}(b) gives the sandbox clause of (ii), and \Cref{prop:twotier:sandbox}(a) shows that sandbox activity cannot affect (i).

\emph{Step 5: depth schedule.} For $B\subseteq\Sigma^P$ let $m_B$ be the least size of a refutation of $B$ ($\infty$ if none) and, for the prefer-unblocked rule, $u_B$ the least size of a refutation with an unblocked descent. Under the default rule $\Refut_d(B)$ is NONE for $d<m_B$ and the lexicographically least refutation of size $m_B$ for $d\ge m_B$: it changes at most once. Under the prefer-unblocked rule it is NONE for $d<m_B$, the least smallest refutation for $m_B\le d<u_B$ (all of them blocked), and the least smallest unblocked refutation for $d\ge u_B$: it changes at most twice. The audit at depth $d$ is determined by the answers $\Refut_d(B)$, $B\subseteq\Sigma^P$, and by $\bigcup\Ccal_d(B_0)$, which is determined by $\Conf_d=\{B:\Refut_d(B)\neq\text{NONE}\}$. There are $2^{|\Sigma^P|}$ sets $B$, each with at most two change points, so the audit's output is piecewise constant in $d$ with finitely many changes and constant for $d\ge d_1^*$. At each stage, Steps 2 and 3 apply with the current $d$; truth maintenance ensures that cached lemmas and macros are derivations in the current $R_A$. For $d\ge d_1^*$ the cleanness of each $\Sigma^*\cup\{\tau\}$ no longer changes, so $\Fres d=F_{\mathrm{res}}$, and the tier is sound relative to $\Rstar\cup R_{F_{\mathrm{res}}}$.

\emph{Non-monotonicity} \status{computed}. Take the practice $\{\mathrm{MP},\mathrm{AC}\}$, no world, and the positions $\pos{A_1}{\emptyset}$ and $\pos{r,\ s\to r}{s}$ with $r,s$ formulas of size~$7$. Then $A=\{\mathrm{MP},\mathrm{AC}\}$ for $d<9$; $A=\emptyset$ for $9\le d<29$, by fallback on the conflict $\{\mathrm{MP},\mathrm{AC}\}$; and $A=\{\mathrm{MP}\}$ for $d\ge29$, where the new conflict $\{\mathrm{AC}\}$ (of size $7+15+7$ on the second position) makes the old one non-minimal and restores MP. Restored schemas are genuine or residual at the larger depth, so soundness is unaffected.

\emph{Uniformity (v).} $G$ is an event about the human samples. On $G$, the accepted set at each time is fixed ($\emptyset$ before $N_1$, $R_A$ after, at each depth of the schedule), whatever the prover and the red team do; so (i)--(iv) hold simultaneously for all of them. \qed

\begin{proof}[Proof of \cref{cor:twotier:truth}]
On $G$, every step the reasoner chains lies in $R_A\subseteq\Rstar\cup R_{\Fres d}$, in $\Ttrust$, or (for the extended reasoner) is a world-channel step. The first two kinds preserve truth in $\mathfrak M$ by hypothesis. A world-channel step has all its judgments in $J_W$ and is accepted only if it is not $W$-falsified; if its premises are true in $\mathfrak M$, they are $W$-true, so its conclusion is not $W$-false, hence $W$-true, hence true in $\mathfrak M$. By induction along derivations, everything derived from $\mathfrak M$-true premises is $\mathfrak M$-true.
\end{proof}

% ---------------------------------------------------------------------
\subsection{Proofs of \texorpdfstring{\cref{prop:twotier:positive,prop:twotier:caution,prop:twotier:negative,prop:twotier:structurality}}{the propositions on positive data, caution, negative evidence and structurality}}\label{app:twotier:caution}

\begin{proof}[Proof of \cref{prop:twotier:positive}]
Every subset of $\Sigma_2$ is clean at every depth, so $\Sigma_1$ and $\Sigma_2$ both satisfy (WS) (\Cref{lem:twotier:onesided}(d)), and the oracle answers are the same under both (\Cref{lem:twotier:onesided}(c)). The learner's view consists of these answers, the fixed prover's queries and its own coins, so it has the same law under both targets. Let the prover query some $q\in R_{\Sigma_2}\setminus\Sound(R_{\Sigma_1})$ in every round. If the learner ever accepts $q$ with probability above $\delta$, it is not $\delta$-sound for $\Sigma_1$; otherwise, under $\Sigma_2$, it fails with probability at least $1-\delta$ to accept the genuine step $q$. In the example, $\vee\mathrm I_1$ is not derivable from the conjunction rules, while $A\wedge B\step B\wedge A$ is derivable by $\wedge$E$_2$, $\wedge$E$_1$ and $\wedge$I.
\end{proof}

\begin{proof}[Proof of \cref{prop:twotier:caution}(a)]
The prover is fixed. The learner's inputs (human data, oracle answers, $W$, escalation answers) have the same law under every $\Sigma\in\mathcal T$: the data law and the practice are shared, oracle answers by \Cref{lem:twotier:onesided}(c), and escalation answers follow the shared practice. Coupling the learner's coins, the whole run up to and including the answer to $q$ has the same law, so $p(h,q)$ does not depend on $\Sigma$. If $q\notin\Sound(R_\Sigma\cup R_{\Fres d(\Sigma)})$ for some admissible $\Sigma$, the event ``$h$, then $q$ accepted'' is an unsound acceptance under $\Sigma$.

\emph{TTL passes.} On $G$ the audit runs on $\hat P=\Sigma^P$, identically in all admissible scenarios. Fix $\Sigma\in\mathcal T$, with valuations $V_\Sigma$ from (WS). If a descent on $\pi$ outputs $s$, then $e^+_\pi$ agrees with $V_\Sigma$, since the agreement argument of \Cref{lem:twotier:descent} needs only a total valuation that extends $e$ and falsifies no trusted step. So $s$ has $V_\Sigma$-true premises and a $V_\Sigma$-false conclusion, and $s\notin R_\Sigma$; hence no schema of $\Sigma$ is removed by descent, and $\Sigma$ is contained in the returned set (descent path) or in $B_0$ (fallback). Let $x\in A\setminus\Sigma$. On the descent path $\Sigma\cup\{x\}\subseteq A$ is $d$-clean by (F1). In the fallback, if $\Sigma\cup\{x\}\subseteq B_0$ were not $d$-clean it would contain some $C\in\Ccal_d(B_0)$, with $C\not\subseteq\Sigma$ because $\Sigma$ is clean, so $x\in C$, contradicting $x\in A$. Either way $x\in\Fres d(\Sigma)$, so $R_A\subseteq R_\Sigma\cup R_{\Fres d(\Sigma)}$.
\end{proof}

\emph{Why joint probabilities and the residue.} The conditional probability given $h$ would be false for randomized learners: a learner that with probability $\varepsilon$ accepts everything from time~1 on, and otherwise runs TTL, is $(\delta+\varepsilon)$-sound, yet given the history ``$q_1$ accepted at time 1'' it accepts every fallacy instance. And plain $\Sound(R_\Sigma)$ would make the intersection trivial, since typically $\emptyset\in\mathcal T$, convicting TTL itself.

\emph{Part (b)} of \Cref{prop:twotier:caution} collects facts proved elsewhere: \Cref{prop:twotier:dilemma}(b); \Cref{thm:coherence:halving}(iii) for the coalition, whose members all contain a falsified fallacy instance under the stated priors (so the coalition-unanimous set does too), while a prior with more than half its mass on $\Rstar$ forces $\Rstar\in S_t$ and hence soundness; and \Cref{prop:search:mdl}.

\begin{proof}[Proof of \cref{prop:twotier:negative}]
Both worlds have the practice $\{\wedge\mathrm I,\mathrm{AC}\}$ with the same frequencies and instance laws, so a learner without refutation evidence has a view with the same law in both. Let the prover query an AC instance outside $\Sound(R_{\wedge\mathrm I})$ under $W_1$, such as $q,\ p\to q\step p$; it is genuine under $W_2$. Accepting it with probability above $\delta$ is unsound under $W_1$; otherwise the learner is incomplete under $W_2$. For the separation, under $W_1$ the instance $(\top,\ \bot\to\top\step\bot)$ has true premises and a false conclusion, while under $W_2$ its premise $\bot\leftrightarrow\top$ is false.
\end{proof}

\begin{proof}[Proof of \cref{prop:twotier:structurality}]
(a) By time $t$ at most $t$ refutations have been exhibited, each a finite derivation implicating finitely many instances; a fallacy with infinitely many invalid instances keeps one that is not refuted. (b) is \Cref{prop:coherence:quus}. (c) With every finite set of steps a hypothesis, the intersection of the version space of data $D$ is $D$ itself (\Cref{sec:imitation:untagged}).
\end{proof}

% ---------------------------------------------------------------------
\subsection{Proof of \texorpdfstring{\cref{thm:twotier:burnin}}{the rare-refuter theorem}}\label{app:twotier:burnin}

\emph{Existence of $s$, and admissibility.} If every $\tau$-instance were in $\Sound(R_{\{\sigma,\mu\}})$, the refutation of $\{\mu,\tau\}$ would transfer, by (F2), to a refutation of $\{\sigma,\mu\}$, which is clean at every depth. So some $\tau$-instance $s$ lies outside $\Sound(R_{\{\sigma,\mu\}})$. The target $\{\sigma,\mu\}$ of scenario~1 is clean at every depth, hence satisfies (WS), and its fallacy $\tau$ has the instance $s$ outside $\Sound(R_{\{\sigma,\mu\}})$; the target $\{\sigma,\tau\}$ of scenario~2 is clean at every depth and has $F=\emptyset$. In scenario~2, $s$ is a genuine instance.

\emph{The inequality.} Let $P_j$ be the law of one human sample (a tag and a step) in scenario $j$. On samples whose tag is not $\mu$, $P_1=(1-\pi)P_2$ as measures: the frequencies are $(1-\pi)q_\sigma$ and $(1-\pi)q_\tau$ against $q_\sigma$ and $q_\tau$, and the per-tag laws coincide. Let $E_t$ be the event that none of the first $t$ samples has tag $\mu$. Then $P_1^{\otimes t}$ restricted to $E_t$ equals $(1-\pi)^tP_2^{\otimes t}$, and $P_2^{\otimes t}(E_t)=1$.

Human samples are i.i.d.\ and independent of everything generated before them. By (OE), every other input of round $r$ is drawn from a kernel applied to the history, the same kernel in both scenarios. By induction over rounds, the joint law of the samples $x=(x_1,\dots,x_t)$ and the other inputs $y$ through time $t$ is $P_j^{\otimes t}(dx)\,K_t(x,dy)$, with one kernel $K_t$ for both $j$. Whether $s$ is accepted at time $t$ is a function of $(x,y)$; let $\phi(x)=K_t(x,\{\text{accept }s\text{ at }t\})\in[0,1]$. Then
\[
\Prob_1(\text{accept }s\text{ at }t)\ \ge\ \int_{E_t}\phi\,dP_1^{\otimes t}\ =\ (1-\pi)^t\int_{E_t}\phi\,dP_2^{\otimes t}\ =\ (1-\pi)^t\,\Prob_2(\text{accept }s\text{ at }t),
\]
so $\Prob_2\le\delta(1-\pi)^{-t}$ whenever $\Prob_1\le\delta$.

\emph{The bounds.} $\Prob_2\ge1-\delta'$ requires $(1-\pi)^t\le\delta/(1-\delta')$, i.e.\ $t\ln\frac1{1-\pi}\ge\ln\frac{1-\delta'}\delta$, which is positive when $\delta+\delta'<1$. Since $\frac1{1-\pi}=1+\frac\pi{1-\pi}$, $\ln\frac1{1-\pi}\le\frac\pi{1-\pi}$, which gives the second inequality. It fails without $\delta+\delta'<1$: for $\delta=\delta'=0.6$ and $\pi=\frac12$ the two sides are $-0.585$ and $-0.405$.

\emph{Tightness.} Let the learner accept $\tau$-instances at time $t$ iff no $\mu$-tag occurs among the first $t$ samples and an independent coin with bias $\min(1,\delta(1-\pi)^{-t})$ comes up heads. It does not escalate, so it satisfies (OE). In scenario~1 it accepts $s$ at time $t$ with probability $(1-\pi)^t\min(1,\delta(1-\pi)^{-t})\le\delta$; in scenario~2, where $\mu$ never occurs, with probability $\min(1,\delta(1-\pi)^{-t})$. Linear programming over all learners that depend only on the tag sequence, for $t\le6$, confirms that this is optimal in all 108 cases checked \status{computed}.

\emph{Practice-following escalation.} By (F2) with $\tau$ and $\mu$ exchanged, some $\mu$-instance $s_\mu$ lies outside $\Sound(R_{\{\sigma,\tau\}})$, hence outside the scenario-2 practice $R_{\{\sigma,\tau\}}$, while it lies in the scenario-1 practice. A learner that escalates $s_\mu$ in round $1$ and accepts $s$ in every later round iff the answer was ``reject'' accepts $s$ with probability $0$ in scenario~1 and $1$ in scenario~2. So without (OE) the conclusion fails; the same linear program, with one practice-following escalation, raises the cap $0$ to $1$ \status{computed}. A lower bound for learners with membership queries to the human would need the refuter to be unknown among many candidates; that is open. \qed

\emph{A concrete instance.} $A_1$ designated, no world, $\sigma=\wedge\mathrm I$, $\tau=\mathrm{AC}$, $\mu=\mathrm{MP}$. The closure of $A_1$ under AC alone is $\{q,\ p\to q,\ p\to\bot,\ p\}$; under $\{\wedge\mathrm I,\mathrm{AC}\}$ it is infinite but contains no $\bot$, because $\wedge$I creates no implications and AC needs an implication with consequent $q$ or $\bot$; $\{\wedge\mathrm I,\mathrm{MP}\}$ is classically sound; and AC followed by MP derives $\bot$.

% ---------------------------------------------------------------------
\subsection{Proof of \texorpdfstring{\cref{thm:twotier:blame}}{the blame theorem}}\label{app:twotier:blame}

(a) The practice $\Sigma_0\cup\{\mu,\tau\}$ and its frequencies are the same in both scenarios, so the data and practice-following escalation answers have the same laws, and oracle answers are target-blind (\Cref{lem:twotier:onesided}(c)). Both targets are clean at every depth, hence satisfy (WS), and by (F2) each target's fallacy has an instance outside the closure of the other's calculus; in particular some $\mu$-instance $s$ lies outside $\Sound(R_{\Sigma_0\cup\{\tau\}})$. In scenario~2, $\mu$ is not residual, because $\Sigma_0\cup\{\tau,\mu\}$ contains the conflict $\{\mu,\tau\}$; so $\Fres d=\emptyset$ there, and accepting $s$ is unsound, which $\delta$-soundness caps at probability $\delta$. The fixed prover queries $s$, and the acceptance probability at any given time is the same in scenario~1, where $s$ is genuine. Finally $\{\mu\}$ and $\{\tau\}$ are clean by (F1), so $\{\mu,\tau\}\in\Ccal_d$ and $\mu\in\Coll_d$ in scenario~1.

(b) Since $d\ge d_0$, a subset of $\Sigma^P$ is $d$-clean iff it is clean at every depth.

\emph{Choice of $M$.} Apply \Cref{lem:twotier:reiter}(b) with $U=B_0$ and $\mathcal K=\Conf_d\cap2^{B_0}$, whose minimal members are $\Ccal_d(B_0)$: since $g\in\bigcup\Ccal_d(B_0)$, some maximal $d$-clean $M\subseteq B_0$ omits $g$. $M$ is clean at every depth, so it satisfies (WS) by \Cref{lem:twotier:onesided}(d) ($e$ is well defined because (WS) holds for the true target).

\emph{Admissibility of $M$.} Let $x\in\Sigma^P\setminus M$. If $x\in B_0\setminus M$, then $M\cup\{x\}$ is not $d$-clean by maximality; if every $x$-instance were in $\Sound(R_M)$, (F2) would turn a refutation of $M\cup\{x\}$ into one of $M$, which is impossible. If $x\in\Sigma^P\setminus B_0$, then $x$ was removed by a descent on some refutation $\pi$, with output $s\in\inst(x)$. The valuation $V_M$ extends $e$ and falsifies no trusted step, so $e^+_\pi$ agrees with it (agreement in \Cref{lem:twotier:descent}); hence $s$ has $V_M$-true premises and a $V_M$-false conclusion, and since $V_M$ preserves truth along $R_M$-derivations, $s\notin\Sound(R_M)$. Either way $x$ has an instance outside $\Sound(R_M)$.

\emph{$\Fres d(M)=\emptyset$.} For $x\in B_0\setminus M$ this is maximality. For $x\notin B_0$, $\{x\}$ is a $d$-conflict by \Cref{lem:twotier:descent}(d), applied in the true scenario, so $M\cup\{x\}$ is not $d$-clean by (F1).

\emph{Witness and conclusion.} $g\in B_0\setminus M$, so the first case gives $s_g\in\inst(g)\setminus\Sound(R_M)$. All admissible scenarios share the data law and the oracle answers, the prover is fixed, and escalation answers follow the shared practice; so the probability that $s_g$ is accepted at a given time is the same under every admissible target. Under $M$, $\delta$-soundness relative to $R_M\cup R_{\Fres d(M)}=R_M$ bounds it by $\delta$. For the schema-level statement, a schema $x$ outside $A=B_0\setminus\bigcup\Ccal_d(B_0)$ lies in $\bigcup\Ccal_d(B_0)$ (handled above) or outside $B_0$, and then its descent output is a witness outside $\Sound(R_M)$ for any maximal clean $M\subseteq B_0$, which is admissible as shown. A policy that includes $x$ in $A'$ accepts that witness when the prover queries it, so includes $x$ with probability at most $\delta$. On the descent path take $M=A$: at stable depth $A$ is clean at every depth, every schema outside $A$ has a descent witness outside $\Sound(R_A)$ by the second case, and $\Fres d(A)=\emptyset$ by \Cref{lem:twotier:descent}(d); so $A$ is admissible and the same argument applies.

(c) \emph{The example.} Take no world, the practice $\{\mathrm{MP},\mathrm{AC}\}$, and $\Ac=\{\pos{A_1}{\emptyset},\pos{A_2}{\emptyset}\}$ with $A_2=\{\varphi,\ \bot\to\varphi\}$, $\varphi=q\wedge(q\wedge q)$. The minimal refutation sizes are $9$ for $\{\mathrm{MP},\mathrm{AC}\}$ (on $A_1$) and $13$ for $\{\mathrm{AC}\}$ (on $A_2$, where AC infers $\bot$ from $\varphi$ and $\bot\to\varphi$), and $\{\mathrm{MP}\}$ is clean at every depth. For $9\le d\le12$ the only minimal conflict is $\{\mathrm{MP},\mathrm{AC}\}$, the descent is blocked, and TTL returns $A=\emptyset$. Yet the admissible targets are only $\{\mathrm{MP}\}$ and $\emptyset$, and asserting $R_{\mathrm{MP}}$ is sound under both (under $\emptyset$, MP is residual). For $d\ge13=d_0$, TTL returns $\{\mathrm{MP}\}$.

\emph{Verification.} These claims were computed with well-founded derivations and size counting distinct judgments \status{computed}, and can be checked by hand. On $A_1$, AC infers $p$ from $q$ and $p\to q$, and MP infers $\bot$ from $p$ and $p\to\bot$; the distinct judgments $q$, $p\to q$, $p$, $p\to\bot$, $\bot$ have $1+3+1+3+1=9$ symbols. On $A_2$, AC infers $\bot$ from $\varphi$ and $\bot\to\varphi$, with $5+7+1=13$ symbols. The closure of $A_1$ under AC alone is $\{q,\ p\to q,\ p\to\bot,\ p\}$, so $\{\mathrm{AC}\}$ has no refutation on $A_1$; and $\{\mathrm{MP}\}$ is classically sound while both contexts are satisfied by $q=1$, $p=0$. So for $9\le d\le12$, $\Ccal_d=\{\{\mathrm{MP},\mathrm{AC}\}\}$. With no world and no trusted steps $e$ is defined only on the assertions and on $\bot$; at $\bot$ the producing MP step has the unevaluated premise $p$ and no premise of value $0$, so the descent is blocked, and the fallback returns $\emptyset$. The subsets of the practice that are clean at every depth are $\emptyset$ and $\{\mathrm{MP}\}$; both are admissible, since AC is classically invalid and MP has instances outside $\Sound(\emptyset)$. Under $\emptyset$, MP is residual at every depth, so asserting $R_{\mathrm{MP}}$ is sound under both. For $d\ge13$, $\{\mathrm{AC}\}$ is a conflict, so $\{\mathrm{MP},\mathrm{AC}\}$ is no longer minimal; with the default oracle the blocked refutation of size $9$ triggers the fallback with $\Ccal_d=\{\{\mathrm{AC}\}\}$, and with the prefer-unblocked oracle the refutation on $A_2$ has an unblocked descent that removes AC; either way $A=\{\mathrm{MP}\}$. \qed

% ---------------------------------------------------------------------
\subsection{Proof of \texorpdfstring{\cref{thm:twotier:depth}}{the depth-relativization theorem}}\label{app:twotier:depth}

\emph{Construction.} Let $\Sigma_0$ be a finite elementary formal system \citep{post1943formal,smullyan1961theory} with judgments $\mathrm{start}_e$ that simulates the universal machine on input $e$ from $\mathrm{start}_e$ and contains the rule $\mathrm{halt}\step\bot$, as in \Cref{thm:coherence:complexity}(a). $\Sigma_0$ has no rules without premises, so without a start judgment nothing is derivable; in particular $\Sigma_0$ is clean on $\pos\emptyset\emptyset$, and the step $\step\mathrm{start}_e$ is not in $\Sound(R_{\Sigma_0})$. Let $\tau_e=(\step\mathrm{start}_e)$ and $\Sigma^P_e=\Sigma_0\cup\{\tau_e\}$, with fixed computable frequencies and instance laws. If $\varphi_e(e)$ diverges, $\Sigma^P_e$ is clean and is the target, with $F=\emptyset$. If $\varphi_e(e)$ converges, $\Sigma_0\cup\{\tau_e\}$ derives $\bot$ from $\emptyset$, the target is $\Sigma_0$, $F=\{\tau_e\}$ and $\tau_e\notin F_{\mathrm{res}}$; accepting $\step\mathrm{start}_e$ is then a step outside $\Rstar\cup R_{F_{\mathrm{res}}}=R_{\Sigma_0}$.

\emph{Uniformity.} For each $e$ only one case occurs, so this is not a two-scenario argument; what is used is that the data laws and the oracle answers $\Refut_d$ (for every $d$, by brute force) are computable uniformly in $e$.

\emph{Reduction.} Let $S=\{e:\exists t\ \Prob(\text{the learner accepts }\step\mathrm{start}_e\text{ by time }t)>\frac12\}$. If $\varphi_e(e)$ diverges, the prover queries the genuine step $\step\mathrm{start}_e$ in every round, and (b) and continuity of measure give $e\in S$. If $\varphi_e(e)$ converges, (a) gives $\Prob(\text{ever accepts})\le\delta<\frac12$, so $e\notin S$. Hence $S$ is the complement of the halting set $\{e:\varphi_e(e)\text{ converges}\}$. But for a computable learner, with computable data laws, oracle and prover, the probability of acceptance by time $t$ is lower semicomputable uniformly in $(e,t)$, so $S$ is $\Sigma_1$. This also holds if the learner escalates, because escalation answers follow the practice by a law computable uniformly in $(e,q)$ (\cref{sec:twotier:lower}), such as $\mathbf 1[q\in R^P_e]$, which is decidable by matching. Computability of the law is needed: $e$ can be read off $\tau_e$, so the law $\mathbf 1[\varphi_e(e)\text{ diverges}]$ depends only on the practice, and a learner that escalates $\step\mathrm{start}_e$ and accepts iff the answer is $1$ would meet (a) and (b). The complement of the halting set is not $\Sigma_1$: a contradiction. \qed

% ---------------------------------------------------------------------
\subsection{Proofs for classical propositional logic}\label{app:twotier:cpc}

\begin{proof}[Proof of \cref{prop:twotier:coherenceonly}]
(a) $\inst(\tau)$ is substitution-closed and contains an unsound step, and $\Rstar$ is complete for $\CPC$, so by \Cref{prop:search:post} $\Cl_{\Rstar\cup\inst(\tau)}(\emptyset)$ contains $\bot$. A derivation of $\bot$ from $\emptyset$ is a refutation of $\Sigma^*\cup\{\tau\}$ relative to $\pos\emptyset\emptyset$, so $d_\tau$ is finite. For $d\ge d_\tau$, $\Sigma^*\cup\{\tau\}\in\Conf_d$ contains some $C'\in\Ccal_d$, and $\tau\in C'$ by \Cref{lem:twotier:onesided}(a). Hence $\Fres d=\emptyset$ for $d\ge\max_\tau d_\tau$, and by \Cref{thm:twotier:main}(iii) and \Cref{cor:twotier:truth} TTL asserts only instances of $\Sigma^*$, which are classically valid.

(b) \emph{$\{\mathrm{AC},K\}$ is a conflict.} $K=A\to(B\to A)$ gives $\bot\to(p\to\bot)$ (with $A=\bot$, $B=p$) and $(p\to\bot)\to(\bot\to(p\to\bot))$ (with $A=p\to\bot$, $B=\bot$). AC, $B,\ A\to B\step A$, applied to these two gives $p\to\bot$; applied to $p\to\bot$ and $\bot\to(p\to\bot)$ it gives $\bot$. This is a refutation relative to $\pos\emptyset\emptyset$. \emph{It is minimal at every depth beyond its size.} $\{K\}\subseteq\Sigma^*$ is clean (\Cref{lem:twotier:onesided}(a)). $\{\mathrm{AC}\}$ is clean at every depth on these positions: from $\emptyset$ it derives nothing, since AC has premises, and the closure of $A_1$ under AC alone is $\{q,\ p\to q,\ p\to\bot,\ p\}$. Likewise, with $A_1$ designated, AC followed by MP derives $\bot$, and $\{\mathrm{MP}\}$ is clean, so $\{\mathrm{AC},\mathrm{MP}\}$ is a minimal conflict. \emph{Every descent is blocked.} With no world and no trusted steps, $e$ is defined only on the assertions of the position and on $\bot$. An output step would need premises among the assertions (or none) and conclusion $\bot$. No genuine step is such a step, because the positions are satisfiable and $\Sigma^*$ is sound; and no AC step is, because AC with premises in $A_1$ yields only $p$, and AC needs premises. So the fallback runs at $B_0=\Sigma^P$, $K$ (and with $A_1$ also MP) lies in $\Coll_d$ for all large $d$, and \Cref{thm:twotier:blame}(b) applies.

(c) $\tau$ derives the denied $\varphi'$ from the asserted $\Pi'$ in one step, a refutation of size at most $|\tau|\le d$ (a closed $\top/\bot$ instance has at most the symbols of $\tau$). So $\{\tau\}\in\Conf_d$ for each $\tau\in F$, every conflict contains a fallacy and hence a singleton conflict, and the minimal conflicts are exactly these singletons. Then $\Coll_d=\emptyset$, $\Fres d=\emptyset$, and \Cref{lem:twotier:audit} gives $A=\Sigma^*$ on the descent path, and also in the fallback, where $\Sigma^*\setminus\Coll_d(B_0)\subseteq A\subseteq\Sigma^*$ and $\Coll_d(B_0)\subseteq\Coll_d=\emptyset$.
\end{proof}

\paragraph{A Hilbert illustration} \status{computed; formulas with at most seven leaves over $\{p,q,\bot\}$, the least bound at which $S$ has instances; with six leaves the conflicts are the same}. Take the practice $\{K,S,\mathrm{MP},\mathrm{DN},\mathrm{AC}\}$ with $\Ac=\{\pos\emptyset\emptyset,\pos{A_1}\emptyset\}$. The minimal conflicts are $\{\mathrm{AC},K\}$ and $\{\mathrm{AC},\mathrm{MP}\}$, the diagnoses are $\{\mathrm{AC}\}$ and $\{K,\mathrm{MP}\}$, and the collateral is $\{K,\mathrm{MP}\}$. The conflict $\{\mathrm{AC},K\}$ is the four-step derivation from no premises in the proof of (b) above. Both conflicts are minimal at every depth beyond their sizes, and with no world and no trusted steps every descent is blocked; so by \Cref{thm:twotier:blame}(b), $K$ and MP are \emph{necessarily} collateral at full depth. With the bilateral position $\pos{A_1}{p}$ the only minimal conflict is $\{\mathrm{AC}\}$, and the closed instance $(\bot\to\bot,\ \bot\to(\bot\to\bot)\step\bot)$ of AC is falsified. Whether $\{S,\mathrm{DN},\mathrm{AC}\}$ is clean at every depth, i.e.\ whether the collateral also contains $S$ or DN, is open. No matrix with at most four values witnesses its cleanness, that is, validates $S$ and DN, with AC designation-preserving, $\bot$ undesignated and $A_1$ satisfiable (exhaustive backtracking, $5.0$M nodes for four values with three designated). A prior that prefers fewer fallacies, Quine's maxim of minimum mutilation, picks $\{\mathrm{AC}\}$ over $\{K,\mathrm{MP}\}$; in the core $\{\mathrm{MP},\mathrm{AC}\}$ case it is silent unless it uses frequency, which does not track validity (\Cref{sec:simplicity}).

% ---------------------------------------------------------------------
\subsection{Proofs for complete decidable theories}\label{app:twotier:decidable}

We take $\Sigma^*$ to be a standard sequent calculus: besides being sound and complete for $T$-consequence, it has $\to$-introduction, $\forall$-introduction, $\neg$-elimination (deriving $\der\bot$ from $\der\psi$ and $\der\neg\psi$), ex falso and weakening, as derived rules at least.

\begin{proof}[Proof of \cref{lem:twotier:complete}]
(a) Let $\Pi\step j$ be a step. \emph{All premises $W$-true.} Each premise $\Gamma\der\varphi$ has a universal closure true in $\mathfrak M$; since $T$ is complete, $T$ proves it, and $\Sigma^*$ derives the premise from $\emptyset$. Hence $\Cl_{\Rstar}(\Pi)=\Cl_{\Rstar}(\emptyset)$, which by soundness and completeness of $\Sigma^*$ and completeness of $T$ is the set of $W$-true sequents; so $j\in\Cl_{\Rstar}(\Pi)$ iff $W(j)=1$. \emph{Some premise $\Gamma\der\varphi$ $W$-false.} By $\to$I and $\forall$I, $\Sigma^*$ derives from it the sentence $\der\forall\bar x(\gamma_1\to\dots\to\gamma_n\to\varphi)$, which is false in $\mathfrak M$. $T$ refutes it, so $\Sigma^*$ derives its negation; then $\neg$E, ex falso and weakening derive every sequent, and $\Cl_{\Rstar}(\Pi)=J$. So $\Pi\step j\in\Sound(\Rstar)$ iff some premise is $W$-false or $j$ is $W$-true, i.e.\ iff the step preserves $W$-truth.

(b) A schema is a fallacy iff it has an instance outside $\Sound(\Rstar)$, iff by (a) it has an instance with $W$-true premises and a $W$-false conclusion. The smallest such instance is a one-step refutation of $\{\tau\}\subseteq\Sigma^*\cup\{\tau\}$ relative to $\pos\emptyset\emptyset$, so $d_\tau\le f_\tau$.
\end{proof}

\begin{proof}[Proof of \cref{cor:twotier:decidable}]
Under the realizability hypothesis \Cref{thm:twotier:main} applies; (WS) holds with $V$ = truth in $\mathfrak M$ (of universal closures), the positions being truthful. The oracle is exact: if $B$ is not $d$-clean it contains a fallacy (\Cref{lem:twotier:onesided}(b)), whose smallest falsified instance has size $f_\tau\le d$ and is found; if $B$ is $d$-clean the search cannot find a falsified instance of size at most $d$, which would be a $d$-refutation. Each $\{\tau\}$, $\tau\in F$, is then a conflict, and the proof of \Cref{cor:twotier:cpc} applies verbatim; truth-soundness follows from \Cref{cor:twotier:truth}, since $\Rstar$ preserves truth in $\mathfrak M$ by \Cref{lem:twotier:complete}(a).
\end{proof}

\paragraph{Why learn at all?} Here $W$ is already a sound and complete step checker, so learning is not needed for epistemic access. TTL adds two things. It identifies \emph{which} calculus the practice uses. And it makes deployed checking fast: schema matching takes linear time, while real quantifier elimination is doubly exponential in the worst case \citep{davenport1988real}, deciding RCF sentences is possible in exponential space \citep{benor1986complexity}, the first-order theory of $(\R,+)$ already needs exponential time, and Presburger arithmetic needs $2^{2^{\Omega(n)}}$ time \citep{fischer1974super}. The audit calls $W$ only on the few instances needed to falsify each fallacy. This dividend is real but modest.

\paragraph{A restricted oracle.} A restricted oracle shows the division of labour more sharply. Suppose $W$ evaluates only quantifier-free sentences with rational constants, and first-order logic is trusted. The fallacy $\tau=(\step\neg\exists x\,(x\cdot x=1+1))$ is a step without premises whose conclusion is not quantifier-free, so it has no instance that $W$ can evaluate: the world alone never refutes it, and \Cref{lem:twotier:complete}(b) fails for this oracle. The position $\pos{\exists x\,(x\cdot x=1+1)}{\emptyset}$ is truthful in $\R$. In the refutation ``$\der\neg\exists x(x\cdot x=1+1)$ (by $\tau$); $\der\bot$ (by trusted $\neg$E, with the asserted $\exists x(x\cdot x=1+1)$)'', the backward rule through $\neg$E gives $\neg\exists x(\dots)$ the value $0$, and the descent outputs $\tau$'s step: $\{\tau\}$ is a singleton conflict. Without trusted logic the conflict is the bag $\{\tau,\neg\mathrm E\}$; with the bilateral position $\pos\emptyset{\der\neg\exists x(x\cdot x=1+1)}$ it is $\{\tau\}$, by a one-step refutation. The position must assert the closed sentence: under the universal-closure semantics, $\pos{x\cdot x=1+1}{\emptyset}$ would assert the false $\forall x\,(x\cdot x=2)$ and violate (WS). Over $\R$, then, coherence with a designated truth catches what this computation misses. Over $\N$ there is no converse: $\Delta_0$ computation refutes nothing that coherence from $\RobQ$ does not already refute (\Cref{lem:coherence:delta0}); what it adds there is localization (\Cref{thm:twotier:arith}(a)).

% ---------------------------------------------------------------------
\subsection{Proof of \texorpdfstring{\cref{thm:twotier:arith}}{the arithmetic theorem}}\label{app:twotier:arith}

(a) In the refutation, $\tau$'s step has no premises and $\forall$E is trusted. The conclusion $\theta(\bar n)$ is a false $\Delta_0$ sentence, so it has value $0$; the backward rule through $\forall$E (whose only premise is $\forall\bar x\theta$; for $k$ variables, through each of the $k$ instantiation steps) gives $\forall\bar x\theta$ the value $0$; and the descent outputs $\tau$'s step, whose (empty) premise set has value $1$ vacuously. The descent is unblocked, so by \Cref{lem:twotier:descent}(c) only schemas having $\der\forall\bar x\theta$ as an instance are removed, and these are fallacies; no axiom of $\PA$ is one, since its instances are true. With the prefer-unblocked oracle, as long as $\tau$ remains in $B$ there is an unblocked refutation of size at most $d$, so the oracle returns an unblocked one, every call removes at least one fallacy and no genuine schema, and the loop does not fall back before $\tau$ is removed. \emph{Size.} For one variable, $|\forall x\theta|=|\theta|+2$, and $\theta(\underline n)$ replaces each of the $m$ occurrences of $x$ (one symbol) by the numeral $S^n0$ ($n+1$ symbols), so $|\theta(\underline n)|=|\theta|+mn$. The total is $2|\theta|+mn+2\le2|\theta|+m(n+1)+O(1)$. For $k$ variables the refutation has $k+1$ judgments, each of size at most $|\theta|+m(n+1)+O(k)$. For $\theta(x)=x<S^50$ ($|\theta|=8$, $m=1$, $n=5$) the sizes are $10$ and $13$, total $23$. Under the default tie-breaking rule a smaller \emph{blocked} refutation through learned axioms could trigger the fallback, and those axioms could become collateral; and if the refutation ran through $\RobQ$'s axioms by coherence instead of through $W$, the bag would contain them. So computation is subsumed by coherence as a source of refutations (\Cref{lem:coherence:delta0}), but not as a source of blame.

(b) $\PA\nder\Con(\PA)$ by G\"odel's second incompleteness theorem \citep{godel1931formal}, so $\PA+\neg\Con(\PA)$ is consistent. A refutation of $\Sigma^*\cup\{\tau\}$ relative to a position with $\Delta_0$ assertions and denials would derive, in first-order logic, from $\PA+\neg\Con(\PA)$ and true $\Delta_0$ leaves, either $\bot$ or a false $\Delta_0$ sentence. $\RobQ\subseteq\PA$ proves every true $\Delta_0$ sentence and refutes every false one (\Cref{lem:coherence:delta0}), so in either case $\PA+\neg\Con(\PA)$ would be inconsistent. Hence $\Sigma^*\cup\{\tau\}$ is clean at every depth and $\tau\in F_{\mathrm{res}}$. Since $\Sigma^P=\Sigma^*\cup\{\tau\}$ is clean at every depth, the first call $\Refut_d(\Sigma^P)$ returns NONE at every $d$, so $A=\Sigma^P\ni\tau$ at every depth (given $\hat P=\Sigma^P$).

(c) \emph{Precise form.} Fix a computable prenex normal form with negations pushed into the $\Delta_0$ matrix, and assert an instance of a practice axiom with $\Sigma_1$ normal form $\exists\bar x\,\theta$ only once $W$ has verified a witness. Then false axioms with $\Sigma_1$ normal form, $\neg\Con(\PA)$ included, are never asserted, and true ones are asserted after a witness search, with no uniform time bound; for schematic axioms, acceptance becomes a semi-decision. The sentence $\forall y\,\exists p\,(\Prf_{\PA}(p,\ulcorner\bot\urcorner)\vee y\neq y)$, equivalent to $\neg\Con(\PA)$ but with a $\Pi_2$ normal form, escapes the policy and is residual.

\emph{Proof.} A false sentence with $\Sigma_1$ normal form $\exists\bar x\,\theta$ has no witness, so $W$ never verifies one and the instance is never asserted. A true one has a witness, which enumeration finds after finitely many $W$-evaluations, with no computable bound on its size. The disguise $\forall y\,\exists p\,(\Prf_{\PA}(p,\ulcorner\bot\urcorner)\vee y\neq y)$ is logically equivalent to $\exists p\,\Prf_{\PA}(p,\ulcorner\bot\urcorner)$, i.e.\ to $\neg\Con(\PA)$, because $y\neq y$ is refutable; its prenex normal form is $\Pi_2$, so the policy does not apply; and $\PA$ plus it is consistent and $\Delta_0$-sound, so by the argument of (b) it is residual.

(d) \emph{Cleanness.} Suppose $\Sigma^*\supseteq\RobQ$ and $\Sigma^*\cup\{\tau\}\cup A\cup\neg D$ is consistent for every designated $\pos AD$. A refutation relative to $\pos AD$ would derive, from $\Sigma^*\cup\{\tau\}\cup A$ and true $\Delta_0$ leaves, either $\bot$, or some $\delta\in D$, or a false $\Delta_0$ sentence. $\RobQ$ proves the true $\Delta_0$ leaves and refutes false $\Delta_0$ sentences, so in each case $\Sigma^*\cup\{\tau\}\cup A\cup\neg D$ would be inconsistent. \emph{Impossibility.} Let $(\varphi_e)$ enumerate the $\Sigma_2$ sentences, let $\Sigma^*=\PA$, and consider the practices $\PA\cup\{\der\varphi_e\}$, with the practice data given as a computable text. Suppose a computable learner, using the data, $\Refut_d$ at every depth and the $\Delta_0$ oracle, eventually and permanently asserts $\der\varphi_e$ if $\varphi_e$ is true, and eventually and permanently withholds it if $\varphi_e$ is false and $\PA\cup\{\der\varphi_e\}$ is clean at every depth. Define $g(e,s)$: if a refutation of $\PA\cup\{\der\varphi_e\}$ has been found by stage $s$ of a systematic search (computable, since each $\Refut_d$ is), output ``false''; otherwise output the learner's current verdict. A true $\varphi_e$ makes the practice sound, hence clean (\Cref{lem:twotier:onesided}(a)), so the search never succeeds and $g$ converges to ``true''. A false $\varphi_e$ either has a refutation, which the search eventually finds, or is clean, and then the learner converges to ``false''. So $\lim_sg(e,s)$ exists and is the truth value of $\varphi_e$, which makes $\Sigma_2$-truth $\Delta_2$ by Shoenfield's limit lemma \citep{shoenfield1959degrees}. But the set of true $\Sigma_2$ sentences is $\Sigma_2$-complete and therefore not $\Pi_2$, let alone $\Delta_2$, by Post's hierarchy theorem (\Cref{thm:coherence:popper}(e)). \emph{Randomized learners} \status{proof sketch}. If the learner is randomized and correct in the limit with probability above $\frac12$, the sets $\{e:\Prob(\text{eventually and permanently withholds }\der\varphi_e)>\frac12\}$ and the corresponding set for assertion are $\Sigma_2$, and the same combination would make $\Sigma_2$-truth $\Delta_2$.

(e) \emph{Precise form.} The induction schema is not a first-order pattern: its lgg has false instances, and the audit refutes it with singleton blame, so the asserted calculus would lack induction. Add auxiliary decidable judgments $\Sub(\varphi,x,t,\psi)$ to $J_W$ and write induction as $\Sub(\varphi,x,0,a),\ \Sub(\varphi,x,Sx,b)\step\ \der a\wedge\forall x(\varphi\to b)\to\forall x\varphi$, with the human data recording the $\Sub$ premises. This is a first-order pattern, recovered by the lgg of generic instances. Then under (Floor) a practice has at most $\min(K,\lfloor1/(2\Delta)\rfloor)$ axiom schemas, and $T_k=\PA+\Con(\PA)+\dots+\Con(T_{k-1})$ is identified from positive data for every $k$ within that bound. Without a floor there are no uniform sample bounds: comparing the targets $T_{k-1}$ and $T_k$, where the tag $\Con(T_{k-1})$ has frequency $\pi$, under (OE), a learner that accepts $\der\Con(T_{k-1})$ at time $t$ with probability at most $\delta$ under $T_{k-1}$ accepts it under $T_k$ with probability at most $1-(1-\delta)(1-\pi)^t$; so if $\delta+\delta'<1$, acceptance with probability at least $1-\delta'$ needs $t\ge\ln\frac{1-\delta}{\delta'}\big/\ln\frac1{1-\pi}$, unbounded as $\pi\to0$. This bound is tight.

\emph{Tag count.} For every tag, $\Delta\le(\rho_i\beta_i-\bar\alpha_i)/2\le\rho_i\beta_i/2\le\pi_i/2$, since $\rho_i\le1$ and $\beta_i\le\pi_i$. As $\sum_i\pi_i=1$, there are at most $1/(2\Delta)$ tags. \emph{Identification.} The axioms of $\RobQ$ and each $\Con(T_j)$ are ground, and fall under the convention $\rho_i=c_i=1$; the $\Sub$-encoded induction rule is a first-order pattern whose lgg on generic instances recovers it: the lgg of five generic $\Sub$-encoded instances equals the pattern up to renaming \status{computed}. So \Cref{thm:imitation:iidnoise} applies and \Cref{thm:twotier:main} gives $\hat P=\Sigma^P$. All practice schemas are genuine (every $\Con(T_j)$ is true), so $F=\emptyset$ and the audit, which never removes genuine schemas, returns $A=\Sigma^P$. Instances of the encoded rule with a false $\Sub$ premise are harmless: $\Sub$ judgments are concluded by no practice schema, so they enter derivations only as leaves, which must be $W$-true. \emph{Without the induction encoding}, the lgg of induction instances is over-general and has false instances: the lgg of four instances has the false instance $(0+0=0)\wedge\forall x(0+0=x\to0+S0=Sx)\to\forall x(0+0=x)$ \status{computed}. For this one, $\forall x(0+0=x\to0+S0=Sx)$ is derivable by trusted equality logic from the true $\Delta_0$ leaf $0+S0=S(0+0)$, so backward propagation through the trusted $\to$E and $\forall$E gives an unblocked descent, and the audit removes the identified schema with singleton blame.

\emph{The floor bound.} Accepting $\der\Con(T_{k-1})$ is unsound under $T_{k-1}$: by G\"odel's second theorem $\Con(T_{k-1})$ is not a theorem of the consistent r.e.\ theory $T_{k-1}$, so the step is not derivable in the target calculus with trusted logic, even with all true $\Sub$ judgments as premises, and it lies outside $\Rstar\cup R_{F_{\mathrm{res}}}=R_{T_{k-1}}$. (It is true, so this is soundness for the target calculus, \Cref{thm:twotier:main}(i), not truth-soundness.) Let $E_t$ be the event that the tag $\Con(T_{k-1})$ does not occur among the first $t$ samples. As in \Cref{app:twotier:burnin}, $P_k^{\otimes t}$ restricted to $E_t$ is $(1-\pi)^tP_{k-1}^{\otimes t}$, $P_{k-1}^{\otimes t}(E_t)=1$, and under (OE) the acceptance probability given the samples is the same function $\phi$ in both scenarios. So
\[
\Prob_k(\text{accept})=\int_{E_t}\phi\,dP_k^{\otimes t}+\int_{E_t^c}\phi\,dP_k^{\otimes t}\ \le\ (1-\pi)^t\Prob_{k-1}(\text{accept})+1-(1-\pi)^t\ \le\ 1-(1-\delta)(1-\pi)^t .
\]
Requiring $\Prob_k\ge1-\delta'$ gives $(1-\pi)^t\le\delta'/(1-\delta)$, i.e.\ $t\ge\ln\frac{1-\delta}{\delta'}\big/\ln\frac1{1-\pi}$, positive when $\delta+\delta'<1$. \emph{Tightness:} the learner that, by one coin at time $0$, accepts the step from the start with probability $\delta$ and otherwise accepts it once its tag has occurred is uniformly $\delta$-sound under $T_{k-1}$, where the tag never occurs, and under $T_k$ accepts at time $t$ with probability $\delta+(1-\delta)(1-(1-\pi)^t)=1-(1-\delta)(1-\pi)^t$. Linear programming over all tag sequences confirms that the bound is attained, in all 108 cases checked \status{computed}. The bound of \Cref{thm:twotier:burnin} does not transfer here: for $\pi=\frac12$, $\delta=10^{-6}$ and $\delta'=0.1$ it would demand $t\ge19.8$, while the learner ``accept iff the tag has occurred'' is $0$-sound under $T_{k-1}$ and accepts under $T_k$ at $t=4$ with probability $0.9375$ \status{computed}. The reason is that the rare tag now lives in the scenario where acceptance is \emph{sound}: the learner must see the tag at least once, so the bound scales with $\ln(1/\delta')$, not $\ln(1/\delta)$. Over classes containing $T_\omega$, identification in the limit fails altogether (\Cref{thm:coherence:turing}); $T_\omega$ has infinitely many axioms, so it lies outside the finite practices considered here.

(f) With $\pos{\Con(\PA)}{\emptyset}$ designated, $\tau=(\step\neg\Con(\PA))$ and the trusted $\neg$E derive $\bot$ from the asserted $\Con(\PA)$. The backward rule gives $\neg\Con(\PA)$ the value $0$, and the descent outputs $\tau$'s step, so $\{\tau\}$ is a singleton conflict (compare \Cref{thm:coherence:residue}(iii): designating true sentences removes the unsound alternatives). \qed

% ---------------------------------------------------------------------
\subsection{Check scripts}\label{app:twotier:scripts}

The figures marked \status{computed} in \cref{sec:twotier} and in this appendix come from the scripts in \texttt{research/theory/T7-checks/} of the accompanying repository, all re-run for this paper. \Cref{tab:twotier:scripts} maps each script to the claims it checks.

\begin{table}[ht]
\centering\small
\begin{tabularx}{\textwidth}{@{}lX@{}}
\toprule
script & claims checked\\
\midrule
\texttt{duality.py} & \cref{lem:twotier:reiter} on 4000 random families\\
\texttt{overlap.py} & strictness of the instance-level inclusion in \cref{prop:twotier:dilemma}(c) (MP/AC common instance)\\
\texttt{noisefree.py} & the noise-free budget $e_i=0$ (\cref{app:twotier:main}, Step~1)\\
\texttt{prop32\_voting.py} & \cref{prop:twotier:voting}: per-position loops versus one query per position\\
\texttt{depth\_nonmono.py} & non-monotonicity in \cref{thm:twotier:main}(iv)\\
\texttt{two\_point\_bounds.py} & tightness of \cref{thm:twotier:burnin}, its failure under practice-following escalation, and the floor bound of \cref{thm:twotier:arith}(e)\\
\texttt{thm56\_finite\_d.py} & the finite-depth counterexample of \cref{thm:twotier:blame}(c)\\
\texttt{hilbert\_blame.py} & the $\to$/$\leftarrow$ witness and the Hilbert illustration (\cref{app:twotier:cpc})\\
\texttt{matrix\_witness.py}, \texttt{matrix4.py} & search for a matrix witnessing that $\{S,\mathrm{DN},\mathrm{AC}\}$ is clean at every depth\\
\texttt{lemma61.py} & \cref{lem:twotier:post} on 30{,}000 random pure schemas\\
\texttt{ind\_lgg.py} & over-general lggs of $\forall$E and of induction (\cref{cor:twotier:decidable}, \cref{thm:twotier:arith}(e))\\
\texttt{ind\_sub\_lgg.py} & the $\Sub$-encoded induction rule is recovered by lgg (\cref{thm:twotier:arith}(e))\\
\texttt{ttl\_sim.py} & \cref{tab:twotier:sim}\\
\bottomrule
\end{tabularx}
\caption{Check scripts for \cref{sec:twotier}.}
\label{tab:twotier:scripts}
\end{table}
